\begin{table}[t]
\begin{tcolorbox}[
sharp corners,
      colframe=refgreen,
      colback=refgreen!5,
      colbacktitle=refgreen!15,
      coltitle=refgreen,
      boxrule=0.4pt,
      title={\textbf{Host-Kernel Decomposition}},
      fonttitle=\bfseries\small,
      top=-10pt, bottom=0pt, left=1pt, right=1pt
]
\begin{tabular}{p{0.48\textwidth} p{0.48\textwidth}}
\begin{lstlisting}[basicstyle=\ttfamily\scriptsize]
# Before decomposing
.../{op_name}/answer/
+-- 0/
    |-- op_host/
    |   |-- add_custom.h
    |   `-- add_custom.cpp
    `-- op_kernel/
        |-- add_custom.h
        `-- add_custom.cpp
\end{lstlisting}
&
\begin{lstlisting}[basicstyle=\ttfamily\scriptsize]
# After decomposing
.../{op_name}/answer/
+-- 0/
|   +-- ... (Original backup)
+-- 1/
|   +-- op_host/ (<- The Host code of scenario 1.)
|   |   +-- add_custom.h
|   |   `-- add_custom.cpp
|   +-- op_kernel/
|   |   +-- add_custom.h
|   |   `-- add_custom.cpp
|   `-- tiling_key.json
`-- 2/
    +-- op_host/ (<- The Host code of scenario 2.)
    |   +-- add_custom.h
    |   `-- add_custom.cpp
    +-- op_kernel/
    |   +-- add_custom.h
    |   `-- add_custom.cpp
    `-- tiling_key.json
\end{lstlisting}
\end{tabular}
\end{tcolorbox}
\caption{The kernel code structure before and after decoupling. }
\label{tab:kernel_host_decompos}
\end{table}